\cotefile{157-2}
\scanpage{291}
\mspage{32}
\uline{Proposition} 7 Soit $f\colon I_{\unsure{\tau|t}}\longrightarrow I_{\tau'}$, un morphisme
d'intervalles \inserted{orientés.} \deleted{\unsure{Soit} \unsure{$f$|$s$} \unsure{une classe}} \deleted{Alors $f$ transforme sommet critique
en sommet critique \unsure{i.e.}}
\[
\deleted{f(I_{\tau\,\unsure{\mathrm{red}|\mathrm{r\acute ed}}})\supset(I_{\tau'})_{\mathrm{red}}}
\]

Soit
\[
 f\colon I_\tau\longrightarrow I_{\tau'}
\]
un morphisme d'intervalles orientés. \deleted{Les}
Considérons les sous-catégories pleines, \unsure{isom.}
à des $\unsure{\mathbb{\Delta}|\Delta}_i^+$ ou $\mathbb{\Delta}_i^-$ comme int. orientés,
\deleted{\unsure{qui}} définies par les intervalles critiques de
$I_\tau$. Ils forment une suite \inserted{$\Sigma_\tau$} de sous catégories
de $I_\tau$
\[
 \Sigma_\tau=(\deleted{\unsure{I_{\tau}}}\mathbb{\Delta}_{l_1}^{\varepsilon_1}(\tau)\unsure{\,|,}
 \mathbb{\Delta}_{l_2}^{\varepsilon_2}(\tau)\,\ldots,\mathbb{\Delta}_{l_r}^{\varepsilon_r}(\tau))
\]
dont $I_\tau$ est la somme amalgamée (on reviendra
sur ce point \unsure{plus}), $r$ étant la longueur utile de
$\tau$. On les appelle les \uline{sous}-\deleted{catégories}\uline{intervalles
critiques} de $I_\tau$. Ainsi $f$ se décompose en
une suite de morphismes
\[
 f_i\deleted{\unsure{I_\tau}}\colon\mathbb{\Delta}_{l_i}^{\varepsilon_i}\deleted{\unsure{\in}}
 \longrightarrow I_{\tau'},\qquad \deleted{\unsure{\forall}}i\in[1,r],
\]
tels que a) pour $i\in[1,r-1]$, le but de $f_i$
soit égal à la source de $f_{i+1}$. \deleted{\unsure{Et que}}\deleted{\unsure{soit d'entre}} De
plus, b) $f_1$ transforme source en source, et $f_r$ transforme
but en but. \deleted{Ces conditions} Inversement, quand
on se donne \unsure{des} $f_i$ satisfaisant ces deux conditions,
ils définissent un morphisme d'intervalles $f$
unique. (Et si la condition b) n'est pas \deleted{\unsure{sat}}
exigée, on trouve une description \unsure{des|de} morphismes
de catégories $I_\tau\longrightarrow I_{\tau'}$.)
% archive sheet number: 290
\scanpage{292}
\mspage{33}
On va d'abord regarder un foncteur quelconque
\[
 f\colon \unsure{\mathbb{\Delta}|\Delta}_i^\varepsilon\longrightarrow I_{\unsure{\tau|t}'}
 \qquad(\varepsilon\in\{\pm1\deleted{\unsure{-}}\}).
\]
Supposons pour \deleted{exemple} \inserted{fixer les idées} que $\varepsilon=+1$, i.e. que \unsure{le|la}
\unsure{premier mb|premier int} soit un intervalle \deleted{\unsure{dire}} critique
\unsure{direct}\ldots{} On trouve que $f$ \deleted{applique} se factorise
par un sous-intervalle critique de $I_{\tau'}$,
\unsure{soit} $\mathbb{\Delta}_j^{\varepsilon'}$.
% slanted
\margin{\uline{Lemme} Si $I_{\tau'}$ est un intervalle, \deleted{\unsure{deux|tout}}\unsure{une|toute} partie $J$ de $I_{\tau'}$ telle que \deleted{\unsure{pour}} pour deux éléments de $J$, il existe tjrs une flèche de l'un à l'autre, est \unsure{contenue|contenu} dans un intervalle critique (unique si $\operatorname{Card}J\geq2$)}
Il est donc défini par une application
croissante de $\Delta_i^\varepsilon$ dans $\Delta_j^{\varepsilon'}$ (si $\varepsilon,\varepsilon'$ de $=$ signe),
une application décroissante dans \unsure{le} sens inverse.

Mais si on exige que $f$ soit croissant, alors
il faut que \uline{ou bien} \deleted{\unsure{$\varepsilon=\varepsilon'$, ou}} $f$ soit \unsure{constante|constant}
(auquel cas \deleted{\unsure{$I_{\tau'}$}}, si sa valeur n'est pas une extrémité
de $I_{\tau'}$, il y a exactement \uline{deux} intervalles
critiques de $I_{\tau'}$ \deleted{qui} \deleted{\unsure{se}} par lesquels se factorise $f$),
ou bien $f$ non constante et $\varepsilon=\varepsilon'$ (et en tous cas,
si $f$ non constante, alors l'intervalle critique par
lequel se factorise $f$ est unique). Pour
résumer :

\uline{Prop 7} Soient $I_\tau,I_{\tau'}$ deux intervalles orientés,
\[
 f\colon I_\tau\unsure{\mathbin{-}|\longrightarrow}I_{\tau'}
\]
un foncteur. Alors pour \deleted{tout} intervalle critique
$J\subset I_\tau$ tel que $f\mid J$ soit non constant, il existe
un unique intervalle critique $J'\subset I_{\tau'}$ tel que
$f\mid J$ se factorise par $J'$ i.e. tel que $f(J)\subset J'$.
Si $f$ est croissant, alors $J$ et $J'$ sont de $=$ signe,
% archive sheet number: 291
\scanpage{293}
\mspage{34}
i.e. tous deux directs, ou tous deux rétrogrades.

Ainsi, si on regarde la suite $s_0,s_1,\ldots,s_{\unsure{m|n}}$
des sommets critiques de \unsure{$\tau$|$I$}, alors \deleted{\unsure{$f$ les $f(s_i)$}}
\[
t_0=f(s_0),\quad t_1=f(s_1)\quad\ldots\quad t_m=f(s_m)
\]
forment une suite de sommets de $I_{\unsure{\tau|t}'}$ \deleted{qui} telle
que deux consécutifs soient reliés par une flèche
de $I_{\tau'}$. Donc la sous-catégorie pleine de
$I_{\tau'}$ définie par ces $t_i$ est \unsure{connexe|convexe}, c'est
donc un sous-intervalle. Si $f$ \deleted{\unsure{transforme}}\deleted{\unsure{est un}}\inserted{est}
croissant, ce sous intervalle est $[t_0,t_m]$.
\deleted{morphisme d'inter}
\deleted{En tous les cas} Si $f$ est un morphisme
d'intervalles, il contient tous les sommets
critiques de $I_{\tau'}$. Dans tous les cas, \deleted{\unsure{tout $t_i$}}
\deleted{\unsure{Si, si}}\deleted{sommet critique \unsure{en}}
prenant $t_i,t_j$ $(i<j)$ et appliquant ce qui précède à $f\mid J$, où $J$ est
le sous-intervalle de $I$ d'origine $s_i$, d'extrémité $s_j$, on voit que tout sommet critique
de $I_{\tau'}$ compris entre $t_i$ et $t_j$ est \deleted{\unsure{intégré}}\inserted{\unsure{l'un|un}}
\unsure{des|les} $t_k$ $(i\leq k\leq j)$. Pour résumer :

\uline{Cor 1} Soient $s_i\unsure{<|\leq}s_j$ deux sommets critiques de
\unsure{$I_{\tau'}$|$I_\tau$}. Alors \deleted{\unsure{pour}} tout sommet critique de $I_{\tau'}$
compris entre $f(s_i)=t_i$ et $f(s_j)=t_j$, \unsure{est} \deleted{\unsure{un}}
l'image d'un sommet critique \inserted{$s_k$} compris entre
$s_i$ et $s_j$ (i.e. $i\leq k\leq j$). En particulier,
si $f$ est un morphisme d'intervalles
\[
f(I_{\tau\,\mathrm{red}})\supset I_{\tau'\,\mathrm{red}}
\]
(où $I_{\tau\,\mathrm{red}}$ désigne l'ens des sommets critiques \inserted{de $I_\tau$},
ou aussi le sous-\deleted{catégorie}intervalle réduit qu'il définit).
% archive sheet number: 292
\scanpage{294}
\mspage{35}
On fera attention que même si $f$ est croissant,
on \uline{n'a pas} en général $f(I_{\unsure{\tau|t}\,\mathrm{red}})\subset I_{\tau'\,\mathrm{red}}$,
i.e. $f$ ne transforme pas forcément sommet critique
en sommet critique. Exemple minimal :

\diagram{\begin{tikzcd}[ampersand replacement=\&,row sep=large,column sep=small]
\unsure{\Phi|\varphi}_3^+\colon \& s_0\arrow[r]\arrow[d,mapsto] \& s_1\arrow[d,mapsto] \& s_2\arrow[l]\arrow[r]\arrow[dl,mapsto] \& s_3\arrow[dl,mapsto]\\
\unsure{\mathbb{\Delta}|\Delta}_2^+\colon \& {}\arrow[r] \& {}\arrow[r] \& {} \& {}
\end{tikzcd}}
$\Phi_3^+$ est réduit, les deux
sommets critiques transformés
en le sommet non critique
$1$.

Désignons par $\mathrm{Cr}(\tau)$ l'ens des intervalles critiques de
$I_\tau$, \deleted{par} en corr. biunivoque (respectant l'ordre total)
avec l'ens des \deleted{\unsure{morphismes identités}}\inserted{flèches de transition} de $I_{\tau\,\mathrm{red}}$.
\[
\mathrm{Cr}(\tau)\simeq\mathrm{Fl}_{\mathrm{tr}}(I_{\tau\,\mathrm{red}}),
\]
Soit, \deleted{$\mathrm{Cr}(\tau,f)$} pour un foncteur
\[
 f\colon I_\tau\longrightarrow\deleted{\unsure{\mathrm{red}}} I_{\tau'\deleted{\unsure{r}}},
\]
$\mathrm{Cr}(\deleted{\unsure{\tau,}}f)$ l'ensemble des intervalles critiques de $I_\tau$
sur lesquels $f$ ne soit pas constante. On a
alors une application
\[
\mathrm{Cr}(\deleted{\unsure{\tau,}}f)\longrightarrow \mathrm{Cr}(\tau')\tag*{(*)}
\]
induite par $f$ via la prop. 7, associant
à \deleted{tout} $J\in\mathrm{Cr}(\tau,\unsure{F|f})$ l'unique $J'\in\mathrm{Cr}(\tau')$
tel que $f(J)\subset J'$. On a de plus

\uline{Corollaire 2} a) L'image de $(*)$ \deleted{contient tout intervalle
critique de $I_{\tau'}$ qui est contenu dans \unsure{$[f(s)$}
le segment d'extrémités $f(s),f(t)$ dans $I_{\tau'}$, et aussi
ceux qui contiennent} \inserted{est un segment de} l'ens. totalement ordonné
$\mathrm{Cr}(\tau')$ i.e. \deleted{\unsure{si deux}}\inserted{si deux} intervalles critiques $J',J''$
\inserted{de $I_{\tau'}$} sont dans l'image, il en est de même de tous
% archive sheet number: 293
\scanpage{295}
\mspage{36}
ceux qui sont compris entre eux.

b) Cette image contient \deleted{tous les} aussi
\deleted{\unsure{Comme de plus} \unsure{$f(\Delta$}}
\deleted{\unsure{les}} intervalles critiques de $I_{\unsure{\tau|t}'}$ (il y en a au plus
deux) qui contiennent \deleted{\unsure{à la}} $f(0_\tau)$ (resp $f(1_\tau)$) et \deleted{qui}\inserted{dont}
l'intersection avec le \uline{segment} $[f(0_\tau),f(1_\tau)]$
est de cardinal $\geq2$, ainsi que tous les
\deleted{segments}\inserted{intervalles} critiques contenus dans ledit segment.

c) Si $f$ est un morphisme \deleted{\uline{croissant}} d'intervalles orientés, alors l'application induite
\[
\mathrm{Cr}(f)\colon\mathrm{Cr}(\deleted{\unsure{\tau,}}f)\longrightarrow\mathrm{Cr}(\tau')
\]
est \deleted{croissante et} \deleted{bijective} \uline{surjective}. \deleted{\unsure{Si $f$ est}} De façon
plus \deleted{\unsure{de}} précise, soit $J'\in\mathrm{Cr}(\tau')$, et soit
$J=f^{-1}(J')$. Alors $\mathrm{Cr}(f)^{-1}(\{\deleted{\unsure{I}}J'\})$ \inserted{est} formé des
int. cr. \deleted{de $I_\tau$}\inserted{$\in\mathrm{Cr}(\tau,f)$} \uline{contenus} \uline{dans} $J$. \deleted{et on a \unsure{long}}
\deleted{\unsure{Si}}On a de plus, \uline{si} $f$ \uline{est \unsure{croissante|croissant}}
\[
\operatorname{long}(J')=
\sum_{\substack{\deleted{\unsure{K\in f^{-1}(f)^{-1}(\{J\})}}\\
 K\in\mathrm{Cr}(f)^{-1}(\{J'\})}}\pi_{\unsure{n|u}}(K,f)
\]
où $\pi_n(K,f)$ est la ``\unsure{portée|pointure}'' \inserted{(numérique)} de $f\mid K$, \deleted{\unsure{son}} i.e.
la longueur de l'intervalle $[f(0_K),f(1_K)]$.

(\uline{NB} Si $K$ est un intervalle critique de $I$ tel que
$K\notin\mathrm{Cr}(f)$ i.e. $f\mid K$ est constante, on posera
$\pi(K,f)=0$. Donc la formule précédente se lit aussi
\[
\operatorname{long}(J')=\sum_{\substack{K\in\mathrm{Cr}(I_\tau),K\subset J\\
(\text{i.e. }K\in\mathrm{Cr}(f^{-1}(\unsure{J|J'})))}}\pi(K,f)
\]

\uline{Cor. \unsure{3|2}} : Si $f$ est un morphisme \uline{croissant} d'int.
orientés, alors pour tout $J'\in\mathrm{Cr}(\tau')$, \deleted{\unsure{le}} on a
\[
(1\leq)\ \operatorname{card}\mathrm{Cr}(f)^{-1}(\{\unsure{J|J'}\})\leq\operatorname{long}J'
\]
% archive sheet number: 294
\scanpage{296}
\mspage{37}
l'égalité n'étant obtenue que si les \unsure{portées|pointures} de $f$
sur les intervalles critiques de $I_{\unsure{\tau|t}}$ \deleted{contenus dans}
qui sont au dessus de $J'$ sont \inserted{toutes} égales à $0$ ou $1$
(donc \inserted{\unsure{ssi}} sur les $\unsure{J|I}\in\mathrm{Cr}(f)^{-1}(\{J\})$, elles sont égales à $1$).

\uline{Cor 4}. Supposons de plus $I'$ \unsure{réduite|réduit}. Alors
\[
\mathrm{Cr}(f)\colon\mathrm{Cr}(\tau,f)\longrightarrow\mathrm{Cr}(\tau')
 \bigl(\simeq\mathrm{Fl}_{\mathrm{tr}}(\tau')\bigr)
\]
est \uline{bijective}.

Donc dans ce cas \inserted{$\tau'$ réduit}, on obtient une application
en sens inverse
\diagram{\begin{tikzcd}[ampersand replacement=\&,column sep=small]
\& \mathrm{Cr}(\tau')\arrow[rr,"\mathrm{Cr}^{\unsure{*|+}}(f)",Rightarrow]\arrow[dr,"\mathrm{Cr}(f)^{-1}"'] \& \& \mathrm{Cr}(\tau)\\
\mathrm{Tr}(\tau')\arrow[ur] \& \& \mathrm{Cr}(\tau,f)\arrow[ur,hook,"\text{inclusion}"'] \& {}
\end{tikzcd}},
application qui transforme \deleted{\unsure{la flèche}}\inserted{l'intervalle}
\deleted{transition}\inserted{critique} $J'$ $(=\{i-1,i\})$ en l'\uline{unique}
intervalle critique $J$ de $I_\tau$ tel que
\[
 f(J)=J'.
\]

% slanted
\margin{\deleted{\unsure{Et pour|pour} \deleted{tout} foncteur $f\colon I_\tau\to I_{\tau'}$, \unsure{pas néc.} croissant ni morphisme d'intervalles,}}
\deleted{Mais on peut aussi \inserted{(sous les \unsure{$=$ hyp. sur $f$})} regarder dans $I_\tau$ l'ens. des
\uline{flèches de transition} \deleted{sur lesquelles}\inserted{ou} (i.e. des
sous-intervalles $J$ de longueur $1$ $\unsure{\{i\mp1,i\}|\{i+1,i\}}$, ou
``\uline{intervalles de transition}''), \inserted{$\unsure{f(u)|f(u_i)}\ne\mathrm{id}$}, \unsure{i.e} telles que
$f\mid J$ non constante. Les \deleted{\unsure{formes}} \unsure{d'où} une
application (supposant $\tau'$ réduit)
\[
\mathrm{Tr}(f)\colon\unsure{\mathrm{Ft}|\mathrm{Tr}}(\tau,f)\longrightarrow\mathrm{Tr}(\tau')
\]
\uline{Cor. 5} : L'image de $\mathrm{Tr}(f)$ est un segment de
l'ens. totalement ordonné $\mathrm{Tr}(\tau')$. Il contient
toutes les flèches de transition \deleted{\unsure{dont le}} de $I_{\tau'}$ dont le}
% archive sheet number: 295
\scanpage{297}
\mspage{38}
\deleted{portée est contenue dans $[f(0_{\unsure{\tau|t}}),f(1_\tau)]$, et est égale
à cet ensemble si $f$ est croissante.}

Mais $J\simeq\unsure{\mathbb{\Delta}|\Delta}_i^\varepsilon$, et on a donc un \deleted{application}\inserted{morphisme}
\uline{surjectif}
\[
 f_0\colon\mathbb{\Delta}_i^\varepsilon\longrightarrow\deleted{\unsure{\mathbb{\Delta}_1}}J'\simeq\mathbb{\Delta}_1^\varepsilon
\]
On voit tout de suite qu'il existe alors une
\uline{unique} flèche de \deleted{\unsure{de $f$}} transition $u$ de
$\mathbb{\Delta}_i^\varepsilon$, telle que $f(u)\ne\mathrm{id}$. Donc on trouve
le corollaire suivant, plus précis que le
cor. 4

\uline{Cor 5} Soit $f\colon I_\tau\unsure{\mathbin{-}|\longrightarrow}I_{\tau'}$ un morphisme
croissant d'intervalles orientés, \deleted{et soit}\inserted{avec} $I_{\tau'}$
réduit, et soit $\mathrm{Tr}(\tau,f)$ l'ens des flèches de
transition \inserted{$\underline{u}_i$} de $I_\tau$ telles que $f(\underline{u}_i)\ne\mathrm{id}$ i.e. telles
que \deleted{la \unsure{portée|pointu}} $\pi(u_i,f)=1$. Alors $f$ induit
une application
\[
\mathrm{Tr}(\tau,f)\longrightarrow\mathrm{Tr}(\tau')
\]
qui est bijective, et (bien sûr) \deleted{croissante}
donc un isom. d'ens. ordonnés.

Ainsi l'application $\mathrm{Tr}(\tau')\unsure{\mathbin{-}|\longrightarrow}\mathrm{Cr}(\tau,f)$ de
tantôt se précise par une application bijective
\[
\mathrm{Tr}(f)^*\colon\mathrm{Tr}(\tau')\longrightarrow\mathrm{Tr}(\tau,f),
\]
reliée à $\mathrm{Cr}(f)^*$ plus haut par le triangle
commutatif
\diagram{\begin{tikzcd}[ampersand replacement=\&,column sep=huge,row sep=small]
\& \mathrm{Tr}(\tau,f)\arrow[dd,Rightarrow,"{\unsure{\sim|\mathord{2}}\ \deleted{\unsure{\varepsilon|E}}\alpha}"]\\
\mathrm{Tr}(\tau')\arrow[ur,"{\mathrm{Tr}(f)^*}"]\arrow[dr,"{\mathrm{Cr}(f)^*}"'] \& {}\\
\& \mathrm{Cr}(\tau,f)
\end{tikzcd}},
% archive sheet number: 296
\scanpage{298}
\mspage{39}
l'application $\alpha$ associant à toute $u_i\in\mathrm{Tr}(\unsure{\tau|t},f)$
l'unique intervalle critique qui ``contienne'' $u_i$
(i.e. tel que $u$ soit une flèche dudit intervalle,
ou encore, \unsure{fasse} partie de la ``\unsure{chaîne|chaine} critique''
\unsure{associée}\ldots).

On a mieux encore :

\uline{Proposition 8} Soient $\tau,\tau'\in\mathrm{Ty}$, \deleted{Alors}\inserted{$\tau'$ \uline{réduit}.} Pour \deleted{tout}
morphisme croissant d'int. orientés
\[
 f\colon I_\tau\unsure{\mathbin{-}|\longrightarrow}I_{\tau'}
\]
(on va \deleted{écrire} dire simplement que $f$ est une flèche
de $\underline{\mathrm{Ty}}$, et écrire $f\in\operatorname{Hom}_{\underline{\mathrm{Ty}}}(I_\tau,I_{\tau'})$),
associons lui la partie $\mathrm{Tr}(\tau,f)$ des
flèches de transition de $I_\tau$ non dégénérées par $f$.
Ces flèches forment un ens. \unsure{de} totalement ordonné,
affecté des signes $+$ ou $-$ (caractère direct ou
rétrograde des flèches), et $f$ induit une \inserted{(cor. 5)}
bijection d'ens. tot. ordonnés
\[
\underset{\substack{\cap\\\mathrm{Tr}(\tau)}}{E_f}=\mathrm{Tr}(\tau,f)\simeq\mathrm{Tr}(\tau'),
\]
et $=$ un isom. d'ens. totalement ordonnés
``affectés de signes''. Ainsi, à tout $f$ on associe
une partie $E_f$ de $\mathrm{Tr}(\tau)$, et telle que
son ``type'' soit égal à $\tau'$. Désignons par
\deleted{$\unsure{\mathscr{P}|\mathcal{P}}_{\tau'}(\unsure{\tau})$ la partie de $\mathscr{P}(I_\tau)$}
\deleted{$\mathscr{P}_{\tau'}(\mathrm{Tr}(\tau))$ la partie de}
\[
\mathscr{P}_{\tau'}(\mathrm{Tr}(\tau))\subset\mathscr{P}(\mathrm{Tr}(\tau))
\]
l'ensemble des parties $E$ de $\mathrm{Tr}(\tau)$ dont le type soit $\tau'$.
% archive sheet number: 297
\scanpage{299}
\mspage{40}
Alors l'application \deleted{\unsure{précéd}}
\[
 f\longmapsto\mathrm{Tr}(\deleted{\unsure{f,}}\unsure{\tau|t},f)\colon
 \operatorname{Hom}_{\underline{\mathrm{Ty}}}(\tau,\tau')\longrightarrow
 \unsure{\mathscr{P}|\mathcal{P}}_{\tau'}(\mathrm{Tr}(\deleted{\unsure{t}}\tau))
\]
est \uline{bijective} : se donner un morphisme
$f\colon\tau\unsure{\mathbin{-}|\longrightarrow}\tau'$ dans $\underline{\mathrm{Ty}}$, \uline{revient au $=$} que
de se donner une suite \inserted{str.} croissante de
flèches de transition
\[
 \deleted{\unsure{u}}\underline{u}_{i_1},\underline{u}_{i_2},\ldots\ \underline{u}_{i_{n'}}
 \qquad(n'=\operatorname{long}\tau')
\]
de $I_\tau$ $(\unsure{u_1|v_1|i_1}<i_2<\ldots<i_{n'})$, telle que
la suite des signatures \unsure{correspondante|correspondantes}
\[
 (\varepsilon_{i_1},\varepsilon_{i_2},\ldots,\varepsilon_{i_{n'}})=\tau'
\]
soit égale à $\tau'$.

\uline{Cor. 1} Pour tout $\tau$, on a des bijections
canoniques
\[
\left\{\begin{aligned}
\mathrm{Tr}^+(\tau)&\simeq\operatorname{Hom}_{\underline{\mathrm{Ty}}}(\tau,\Delta^+)\\
\mathrm{Tr}^-(\tau)&\simeq\operatorname{Hom}_{\underline{\mathrm{Ty}}}(\tau,\Delta^-)
\end{aligned}\right.,
\]
\uline{fonctorielles en $\tau$} (variant dans $\underline{\mathrm{Ty}}$).

Pour donner un sens au corollaire, il
faut définir \deleted{$\tau\mapsto\mathrm{Tr}^+(\tau)$ et $\mathrm{Tr}^-(\tau)$ comme}
\[
 \tau\unsure{\mapsto|\mathrel{\vdash}}\mathrm{Tr}^{+}(\tau),\qquad \tau\mapsto\mathrm{Tr}^{-}(\tau)
\]
comme \uline{foncteurs} contravariants en $\tau$, donc
pour un morphisme
\[
 f\colon I_{\tau_{\unsure{2|1|\,}}}\unsure{\mathbin{-}|\longrightarrow}I_{\tau_{\unsure{1|2|\,}}}
\]
dans $\mathrm{Ty}$, définir \unsure{des} applications
% archive sheet number: 298
\scanpage{300}
\mspage{41}
\[
\begin{aligned}
\mathrm{Tr}^+(f)^*&\colon\mathrm{Tr}^+(\unsure{\tau|t}_1)\unsure{\mathbin{-}|\longrightarrow}\mathrm{Tr}^+(\tau_2)\\
\mathrm{Tr}^-(f)^*&\quad\mathrm{Tr}^{\deleted{\unsure{+|-}}\unsure{-|+}}(\tau_1)\mathbin{-}\mathrm{Tr}^-(\tau_{\unsure{1|2}}),
\end{aligned}
\]
en d'autres termes, une application
\[
 \mathrm{Tr}(f)^{\unsure{*|'}}\colon\mathrm{Tr}(\tau_1)\mathbin{-}\mathrm{Tr}(\tau_2)
\]
``compatible avec les signatures''. Dans le cor. 5
plus haut, on a défini une telle application seulement pour $\tau'$ réduit.
% slanted
\margin{Ce qui suit ne fait pas intervenir les signatures -- c'est un \unsure{sorite|sorites} sur les applications croissantes entre \unsure{ens.} finis totalement ordonnés.}
Pour une flèche \deleted{\unsure{de $f$}}
de $I_{\unsure{\tau,|\tau_1}}$, \deleted{désign} appelons \uline{portée ensembliste} de
$f$ sur $u$, ou de $u$ par $f$, la portée \inserted{ensembliste} de
la flèche $f(u)$ dans $I_{\tau'}$, i.e. \deleted{\unsure{le segment}}
le segment $[f(i-1),f(i)]\subset I_{\tau'}$ (si $u=\underline{u}_{\unsure{i}}$).
\deleted{On voit alors que les portées des}

Soient $\underline{u}_{i_1},\ldots,\underline{u}_{i_N}$ ou $(v_1,v_2,\ldots,v_N)$
la suite croissante des éléments de $\mathrm{Tr}(\tau,f)$,
on voit \inserted{\unsure{donc par réc. sur $i$|disons par réc. sur $i$}} que pour \deleted{\unsure{$1<i\leq N$}} $i\in[1,N-1]$, on a
\[
\left\{\begin{aligned}
 f(\operatorname{but}(v_i))&=f(\operatorname{source}(v_{i+1}))\deleted{\unsure{=}}\ ( = t_i)\\
 \deleted{\text{et si on désigne}}\quad f(\operatorname{source}v_1)&=0_{\tau'},\quad f(\operatorname{but}(v_N))=1_{\tau'}^{\unsure{+}},
\end{aligned}\right.
\]
Considérons la suite
\[
\{\overset{\substack{t_0\\\shortparallel}}{0_{\tau'}},t_1,t_2,\ldots\ t_{N-1},\overset{\substack{t_N\\\shortparallel}}{1_{\tau'}}\}
\]
elle est définie par le sous-ensemble correspondant de $I_{\tau'}$, qu'on appellera l'ensemble des
\uline{sommets critiques} \uline{de} $I_{\tau'}$ \uline{pour} $f$. \deleted{Donc} C'est
une suite \inserted{contenant $0_{\tau'},1_{\tau'}$ et} telle que deux éléments consécutifs
sont reliés dans $I_{\tau'}$ par une flèche, ce qui \deleted{implique}
\inserted{équivaut à dire} qu'elle contient les sommets critiques (intrinsèques)
de $I_{\tau'}$. \deleted{\unsure{Or}} D'autre part, pour les \deleted{flèches de}\inserted{intervalles}
de transition $[i-1,i]$ de $I_{\tau'}$ \deleted{(qui} $(i\in[1,n'])$,
% archive sheet number: 299
\scanpage{301}
\mspage{42}
il existe un et un seul intervalle de transition
\deleted{\unsure{pour}} $[t_{j-1},t_j]$ pour cette suite $(j\in[1,N])$, \unsure{telle|tel}
qu'on ait $[i-1,i]\subset[t_{j-1},t_j]$. C'est dire qu'il
existe une et une seule $v_j\deleted{\unsure{\underline{u}_j}}\in\mathrm{Tr}(\unsure{\tau|t},f)$ telle
que la portée de $f$ sur $\deleted{\unsure{\underline{u}_j}}v_j$ (savoir justement
$\unsure{t_{j-1},t_j|[t_{j-1},t_j]}$) contienne celle de $\underline{u}_i$. Pour résumer

\uline{Corollaire 2} Soit $f\colon\tau\longrightarrow\tau'$ dans $\underline{\mathrm{Ty}}$. Alors
pour \deleted{\unsure{tout|toute}} $u\in\mathrm{Tr}(\tau')=\mathrm{Tr}(I_{\tau'})$, il existe une
et une seule $\deleted{\unsure{u}}v\in\mathrm{Tr}(\tau)$ telle que la portée
de $f$ sur $v$ contienne $u$. \deleted{\unsure{et}} De plus, les
signatures de $u$ et de $v$ \unsure{sont} \deleted{les mêmes} égales.

C'est ainsi qu'on obtient donc une
application
\[
\mathrm{Tr}(f)^*\colon\mathrm{Tr}(\tau')\longrightarrow\mathrm{Tr}(\tau\deleted{\unsure{,f}})
\quad\text{(se factorisant par $\mathrm{Tr}(\tau,f)$)}
\]
induisant
\[
\begin{aligned}
\mathrm{Tr}^+(f)^*&\colon\mathrm{Tr}^+(\tau')\longrightarrow\mathrm{Tr}^+(\tau)\\
\mathrm{Tr}^-(f)^*&\colon\mathrm{Tr}^-(\tau')\unsure{\mathbin{-}|\longrightarrow}\mathrm{Tr}^-(\tau).
\end{aligned}
\]

La transitivité pour un composé
\diagram{\begin{tikzcd}[ampersand replacement=\&]
\tau^{\unsure{\leftarrow|\mathbin{-}}}\arrow[r,"f",dash] \& \tau'\arrow[r,"g",dash] \& \tau''
\end{tikzcd}}
ou
\diagram{\begin{tikzcd}[ampersand replacement=\&]
I_\tau\arrow[r,"f",dash] \& I_{\tau'}\arrow[r,"{\unsure{f|g}}",dash] \& I_{\tau''}
\end{tikzcd}}
est immédiate. Si $u,u',u''$ sont des fl. de tr.
dans $I_\tau,I_{\tau'},I_{\tau''}$ respectivement, telles que
\[
\begin{aligned}
\pi(g,u')&\supset\pi(u'') &\text{i.e. } u'&=\mathrm{Tr}(g)^*(u'')\\
\pi(f,u)&\supset\pi(u') &\text{i.e. } u&=\mathrm{Tr}(f)^*(u')
\end{aligned}
\]
alors on a
\[
\pi(gf)(u)\supset\pi(u'')\qquad\text{i.e. }u=\mathrm{Tr}(gf)^*(u'').
\]
% archive sheet number: 300
\scanpage{302}
\mspage{43}
\deleted{Introduisant les intervalles critiques $J,J',J''$ qui
contiennent \deleted{\unsure{J}} $u,u',u''$ respectivement, on est ramené}
\deleted{\unsure{à}} Si les flèches en question ont comme portée
intrinsèque $[i-1,i]$, $[i'-1,i']$, $[i''-1,i'']$ respectivement,
l'hypothèse est
\[
\begin{aligned}
 g(i'-1)&\leq i''-1,\quad g(i')\geq i''\\
 f(i-1)&\leq i'-1,\quad f(i)\geq i'
\end{aligned}
\]
\deleted{\unsure{et}} comme $g$ est croissant, on a d'ailleurs
\[
\begin{aligned}
 gf(i-1)&\leq g(i'-1)\ (\leq i''-1),\\
 gf(\unsure{u|i})&\geq g(i')\ (\geq i''),
\end{aligned}
\qquad\text{cqfd.}
\]

Ainsi, le cor. 1 plus haut obtient un sens précis,
et il est démontré du $=$ coup, vu que les
correspondances du cor. 1 sont données par
\[
\underset{\substack{\rotatebox{-90}{$\in$}\\
\operatorname{Hom}_{\underline{\mathrm{Ty}}}(\unsure{\tau|t},\unsure{\mathbb{\Delta}|\Delta}^{\unsure{\varepsilon|+}})}}{f}
\longmapsto\mathrm{Tr}(f)^*(\underline{u})
\]
où $\underline{u}$ est l'unique flèche de transition de $\mathbb{\Delta}^\varepsilon$.

On va utiliser ces sorites pour donner,
pour toute $\tau\in\mathrm{Ty}$, une description combinatoire
simple des catégories $\underline{\mathrm{Ty}}_{/\tau}$ et $ {}_{\tau\backslash}\underline{\mathrm{Ty}}$, ou du
moins de leurs ensembles d'objets.

\uline{1\textsuperscript{o}} \uline{Description de} $\operatorname{Ob}\,{}_{\tau\backslash}\underline{\mathrm{Ty}}\deleted{\unsure{/\tau}}$.
À $f\in\operatorname{Ob}\,{}_{\tau\backslash}\underline{\mathrm{Ty}}\deleted{\unsure{\tau}}$ i.e.
\[
 f\colon\tau\longrightarrow\tau'\qquad\text{i.e. } f\colon I_\tau\unsure{\mathbin{-}|\longrightarrow}I_{\tau'}
\]
dans $\underline{\mathrm{Ty}}$, est associé $E\colon E_f=\deleted{\unsure{\mathrm{Cr}}}\mathrm{Tr}(\tau,f)$\inserted{$\subset\mathrm{Tr}(\tau)$}, l'ens des flèches
\inserted{de transition} critiques (i.e. non dégénérées) par $f$. À chaque
$u\in\deleted{\unsure{\mathrm{Tr}}}E$, associons l'entier $\operatorname{long}(\pi(f,u))$,
% archive sheet number: 301
\scanpage{303}
\mspage{44}
la \uline{portée} de $f$ sur $u$. Alors \deleted{$\unsure{\tau|E}$} $(\tau',f)$ est déterminé
de façon unique par le couple
\[
 (E,\pi),\quad E\in\unsure{\mathscr{P}|\mathcal{P}}(\mathrm{Tr}(\tau)),\quad
 \underset{\text{\(\hookrightarrow\) pondération}}{\pi}\colon E\longrightarrow
 \underset{\substack{\parallel\\\text{ens. des}\\\text{entiers }\geq 1}}{\mathbb N^*}.
\]
Ou si on préfère, par l'application $\pi=\pi_f$
\[
 \mathrm{Tr}(\tau)\xrightarrow{\pi}\mathbb N,\quad
 u\longmapsto\pi(u)=\pi(u,f)=(\inserted{\operatorname{long}}\pi_{\tau'}(f(u)))
\]
(qui redonne $E$ comme
\[
 E=\pi^{-1}(\mathbb N^*)\;).
\]
Cette application n'est soumise à aucune restriction
\[
 \operatorname{Ob}\,{}_{\tau\backslash}\underline{\mathrm{Ty}}\deleted{\unsure{/\tau}}
 \xrightarrow[\ ]{\overset{\mathrm{bij.}}{\sim}}
 \deleted{\unsure{\mathrm{Appl.}|\mathrm{Appl}}}\operatorname{Hom}_{\mathrm{Ens}}(\mathrm{Tr}(\tau),\mathbb N)
\]
On récupère $\tau'$ à partir de $\pi$ ainsi. On
prend $E=\pi^{-1}(\mathbb N^*)$ (fl. de transition \deleted{\unsure{pour}} \inserted{non dégénérées}
pour $\pi$), $E$ est \deleted{\unsure{totalemt|totalement}} ordonné par l'ordre induit
par celui de $\mathrm{Tr}(\tau)$, il est de plus signaturé
par les signatures de $\mathrm{Tr}(\tau)$, soient
\[
 (\varepsilon_1,n_1)\mathbin{-}(\varepsilon_r,n_r)
\]
la suite des signatures et pondérations, \deleted{associées} \inserted{de la}
suite des éléments de $E$. \quad ($\varepsilon_i\in\{\pm1\}$, $n_i\in\mathbb N^*$
pour $i\in[1,r]$). Alors,
\[
 \tau'=\Delta_{\varepsilon_1}^{\deleted{\unsure{\varepsilon_1}}n_1}\Delta_{\varepsilon_2}^{n_2}
 \deleted{\unsure{\Delta_{\varepsilon_3}^{n_3}}}\ldots\Delta_{\varepsilon_r}^{n_r}
 \qquad r=\operatorname{card}E
\]
dans le monoïde libre $\mathrm{Ty}$, où on a désigné
par $\Delta_+,\Delta_-$ (au lieu de $\Delta^+,\Delta^-$) les
deux générateurs. \unsure{et application|et l'application}
\[
 I_\tau\unsure{\mathbin{-}|\longrightarrow} I_{\tau'}
\]
% archive sheet number: 302
\scanpage{304}
\mspage{45}
% archive sheet number: 303
est décrite de façon unique en disant qu'elle est \uline{constante}
\deleted{en dehors \deleted{de la} des intervalles \unsure{flèches|lignes} de transition}
sur chacune des ``composantes connexes'' (au sens de
la numérotation des sommets \unsure{\}|)} de $I_\tau$) de l'ens
\deleted{$\unsure{E}$}\inserted{$E'\deleted{\unsure{(f)}}=E'(f)\subset I_\tau=[0,n]$}
des sommets qui sont sources ou buts de flèches
de transition dégénérées pour $f$ i.e. de $\mathrm{Tr}(\tau)\setminus E$.
Et sur le $\alpha$.ième élém\textsuperscript{t} de $E$ \inserted{d'indice $i_\alpha$ dans $\mathrm{Tr}(\tau)$,}
\deleted{$(\unsure{1\leq i\leq r})$}
($\alpha\in[1,r=\operatorname{card} E]$), $f$ induit le foncteur
\diagram{\begin{tikzcd}[ampersand replacement=\&,column sep=small,row sep=small]
 \& J_{i_\alpha}\arrow[r]\arrow[d,dash] \& I_{\tau'}\\
 \unsure{\mathbb{\Delta}|\Delta}_{\varepsilon_\alpha}\arrow[ur,phantom,"\simeq" description,sloped] \&
 \substack{\text{intervalle}\\\text{de transition}\\\text{d'indice }i_\alpha\\\text{dans }I_\tau} \&
\end{tikzcd}}
donné par la \deleted{flèche} \uline{unique} \inserted{$\underline{\mathrm{Ty}}$-}\uline{flèche} de $I_\tau$
de $J_{i_\alpha}$ dans le $\alpha$.ième facteur, \inserted{de type} $\Delta_{\varepsilon_\alpha}^{n_\alpha}$, de $\tau'$\unsure{$*$|,}
donné par $(*)$
\[
 \Delta_{\varepsilon_\alpha}^{n_\alpha}\simeq\text{type de }\mathbb{\Delta}_{n_\alpha}^{\varepsilon_\alpha}\quad\ldots
\]

2\textsuperscript{o} \uline{Détermination de} $\operatorname{Ob}\,\underline{\mathrm{Ty}}_{/\tau'}$

Si $(\tau,f)\in\operatorname{Ob}\,\underline{\mathrm{Ty}}_{\unsure{\tau'|/\tau'}}$, \unsure{donc}
\[
 f\colon I_\tau\longrightarrow I_{\tau'},
\]
on lui associe \deleted{\unsure{donc}} d'abord la partie
\[
 \unsure{\mathrm{Cr}|\mathrm{Tr}}(\tau',f)=E'_f\inserted{{}=E'}\subset I_{\tau'}
\]
\inserted{formée} des sommets \inserted{(dits ``critiques pour $f$'', cf page 41)}
de $I_{\tau'}$ \deleted{\unsure{critiques} qui sont séparés,}
\inserted{\unsure{disons en l'un|disons en un} bout d'une flèche $f(u)$, où $u\in\mathrm{Tr}(\tau,f)$,}
% slanted
\deleted{c'est une partie de $I_{\tau'}$ contenant les sommets critiques}
i.e. \inserted{\unsure{$f$ croît strictement|$f$ croisse strictement}} \unsure{$f(u)$|$f u$} n'est pas dégénérée.
\deleted{\unsure{Celle|Cette} \unsure{Celle|Cette}} $E'$
contient \deleted{\unsure{en particulier les} $f$\unsure{en} $+$}
$f(\deleted{\unsure{E}} I_{\tau\,\mathrm{red}})$ (i.e.
les images par $f$ des sommets critiques de $I_\tau$)
\scanpage{305}
\mspage{46}
mais peut-être plus grand. (P. ex. si $\tau=\tau'$, $f=\mathrm{id}$,
\deleted{\unsure{cela contient}} $E'=I_{\tau'}$, qui est donc str. plus grand que $I_{\tau'\,\unsure{\mathrm{r\acute ed}|\mathrm{red}}}$
si $\tau'$ n'est pas réduit). Pour \deleted{tout} $i\in E'$,
considérons $f^{-1}(\{i\})\subset I_\tau$. \inserted{Comme $f$ croissant,} c'est un intervalle
de l'ens \deleted{\unsure{tot.}} ordonné $I_\tau$, donc c'est un sous-intervalle
de \unsure{la|l'} ``\unsure{catégorie|catigo|catég}'' intervalle $I_\tau$. Comme
tel, il a une signature $\tau_i\in\mathrm{Ty}$. On a donc
une application
\[
 \unsure{\eta|\upsilon}_f\deleted{\unsure{=}}\ \text{ou }\eta\colon E'\longrightarrow\mathrm{Ty}.
\]
Je dis que $(\tau,f)\in\operatorname{Ob}\,\underline{\mathrm{Ty}}_{/\tau'}$ est déterminé de
façon unique par ce couple
\[
 (E',\eta)\ \left|\begin{aligned}
 E'&\in\unsure{\mathscr P|\mathcal P}(I_{\tau'})\unsure{=|\colon}\mathscr P([0,n'])\\
 \eta&\colon E'\longrightarrow\mathrm{Ty}
 \end{aligned}\right.\quad,
\]
De plus, celui-ci est soumis à
\deleted{aucune} \inserted{la seule restriction} que \deleted{$E'$}
\[
 E'\supset\mathrm{Cr}(\tau')=\substack{\text{ens. des sommets critiques}\\\text{de }I_{\tau'}\, .}
\]
(ou ce qui revient au $=$, que $E'$ contient $0_{\tau'},1_{\tau'}$
et que deux éléments consécutifs de $E'$ sont
reliés par une flèche de $I_{\tau'}$)\unsure{;|:} \deleted{\unsure{le}}
\[
 \left\{\begin{aligned}
 \operatorname{Ob}\big(\underline{\mathrm{Ty}}_{/\tau'}\big)&\simeq\text{ens des couples }(E',\eta),\ \text{avec}\\
 &E'\subset[0,\underset{\substack{\parallel\\\operatorname{long}\tau'}}{n'}],\quad E'\supset\mathrm{Cr}(\tau')\\
 &\eta\colon E'\longrightarrow\mathrm{Ty}
 \end{aligned}\right.
\]
Soit
\[
 0=t_0<t_1<\ldots<t_N=n'
\]
la suite des éléments de $E'$, soit pour \tikz[baseline=(q.base)]\node[draw,ellipse,inner sep=2pt](q){$i\in[1,N]$};
\[
 \left\Vert\begin{aligned}
 \unsure{l_i|h_i}&=t_i-t_{i-1} &\unsure{\neq1|\geq1}\\
 \varepsilon_i&\quad\text{signature de }[t_{i-1},t_i]\ \text{dans }I_{\tau'}\\
 &\deleted{\unsure{I_{\tau}\ \eta(\ }} 
 \end{aligned}\right.
\]
% archive sheet number: 304
\scanpage{306}
\mspage{47}
et pour \tikz[baseline=(q.base)]\node[draw,ellipse,inner sep=2pt](q){$i\in[0,N]$};, soit
\[
 \left\Vert\,\tau_i=\unsure{\eta|\upsilon}(t_i)\,.\right.
\]
Alors le type $\tau$ est donné par
\[
 \begin{aligned}
 \tau&=\tau_0\Delta_{\varepsilon_1}^{\deleted{\unsure{l_1|h_1}}}\tau_1
 \Delta_{\varepsilon_2}^{\deleted{\unsure{l_2|h_2}}}\ldots
 \Delta_{\varepsilon_N}^{\deleted{\unsure{l_N|h_N}}}\tau_N\\
 &\qquad\text{longueur}\qquad N+\sum_0^N\underset{\substack{\rule{0.3pt}{3pt}\rule{3.6em}{0.3pt}\rule{0.3pt}{3pt}\\\lambda_i\geq0}}{\operatorname{long}(\tau_i)}
 \end{aligned}\tag*{$(*)$}
\]
3\textsuperscript{o}) Description de \deleted{$\unsure{\mathrm{Fl}(\operatorname{Ob}|\operatorname{Ob}}$} $\mathrm{Fl}(\mathrm{Ty})$.

Si
\[
 f\colon\tau\longrightarrow\tau'\qquad\text{dans }\mathrm{Ty}
\]
on lui associe comme tantôt un entier $\uuline{N}\geq0$
($N=\operatorname{card}\mathrm{Cr}(\tau',f)-1$, \deleted{il est} $N=0$ ssi $\tau'=1$,
type de l'intervalle ponctuel), \uline{des} applications
\[
 \left\{\begin{aligned}
 [1,N]&\longrightarrow[+1,-1]\times\mathbb N^*,& i&\longmapsto(\varepsilon_i,\unsure{l_i|h_i})\\
 [0,N]&\longrightarrow\mathrm{Ty},& i&\longmapsto\tau_i\,.
 \end{aligned}\right.
\]
% slanted
\margin{\uline{NB} Si $N=0$ $[1,N]=\varnothing$ et il n'y a pas de $\varepsilon_i,l_i$ ; l'application
$[\underset{\substack{\parallel\\\{0\}}}{0,N}]\unsure{\mathbin{-}|\longrightarrow}\mathrm{Ty}$
se réduit à la donnée de $\tau_0\in\mathrm{Ty}$.}
La flèche $f$ dans $\mathrm{Ty}$ est déterminée par ces
données, le type $\tau$ est donné par la formule
$*$, ci-dessus, le type $\tau'$ par
\[
 \tau'=\Delta_{\varepsilon_1}^{l_1}\Delta_{\varepsilon_2}^{l_2}\ldots
 \Delta_{\varepsilon_N}^{l_N},\qquad\substack{\text{longueur}\\n'=\sum l_i}
 \tag*{$(*')$}
\]
et l'application $f$ croissante de $\inserted{I_\tau={}}[0,N+\sum\operatorname{long}\tau_i]$
dans $I_{\tau'}=[0,\sum l_i]$ est donnée ainsi
\[
 \begin{aligned}
 \deleted{\unsure{\Rightarrow}}\ \text{constante sur } &[0,0+\lambda_0],&\deleted{\unsure{\longmapsto}}\ 0\\
 \deleted{\unsure{\lambda_0+1|0+1}}\ \text{sur }&
 [\lambda_0+1,\lambda_0+\unsure{1|l}\deleted{\unsure{+1}}+\lambda_1]&\longmapsto l_1\\
 \text{\unsure{constant|constante} sur }&[\lambda_0+\lambda_1+\unsure{1|2},\lambda_0+\lambda_1+2+\lambda_2]&\unsure{\mathrel{\vdash\mkern-7mu\dashv}|\longmapsto}l_2\\
 \ldots&\\
 \text{\unsure{constante|constant} sur }&[\lambda_0+\lambda_{\unsure{1|i}}\deleted{\unsure{+}}\ldots+\lambda_i+(i+1),
 \lambda_0+\ldots+\lambda_i+\lambda_{i+1}+\unsure{(i+1)|(i-1)}]\unsure{\mathrel{\vdash\mkern-7mu\dashv}|\longmapsto}l_i
 \end{aligned}
\]
% archive sheet number: 305
